\documentclass[11pt]{article}
\usepackage[margin=1in]{geometry}
\usepackage{amsmath,amssymb,amsthm}
\usepackage{hyperref}

\newtheorem{theorem}{Theorem}
\newtheorem{lemma}[theorem]{Lemma}
\newtheorem{definition}[theorem]{Definition}
\newtheorem{example}[theorem]{Example}

\title{Continuity: $\epsilon$-$\delta$ and Sequential Characterisations}
\author{Math Foundations --- Node 11}
\date{}

\begin{document}
\maketitle

\section{Definitions}

Let $(X, d_X)$ and $(Y, d_Y)$ be metric spaces. We routinely specialise
to $X = Y = \mathbb{R}$ with $d(x, y) = |x - y|$.

\begin{definition}[Continuity at a point]\label{def:cont}
A function $f : X \to Y$ is \emph{continuous at $c \in X$} if for every
$\epsilon > 0$ there exists $\delta > 0$ such that
\[
  d_X(x, c) < \delta \implies d_Y(f(x), f(c)) < \epsilon.
\]
$f$ is \emph{continuous on $X$} if it is continuous at every point.
\end{definition}

\begin{definition}[Uniform continuity]\label{def:unif}
$f$ is \emph{uniformly continuous on $X$} if for every $\epsilon > 0$ there
exists $\delta > 0$ such that for all $x, y \in X$,
\[
  d_X(x, y) < \delta \implies d_Y(f(x), f(y)) < \epsilon.
\]
The single $\delta$ must serve every pair of points; the quantifier order
distinguishes this from pointwise continuity.
\end{definition}

\begin{definition}[Lipschitz continuity]\label{def:lip}
$f$ is \emph{Lipschitz with constant $L \geq 0$} if
\[
  d_Y(f(x), f(y)) \leq L\, d_X(x, y) \quad \text{for all } x, y \in X.
\]
\end{definition}

The trivial chain of implications holds:
\[
  \text{Lipschitz} \implies \text{uniformly continuous} \implies \text{continuous}.
\]
None of the converses hold: $\sqrt{x}$ on $[0,1]$ is uniformly continuous but
not Lipschitz; $x^2$ on $\mathbb{R}$ is continuous but not uniformly continuous.

\section{Main theorem}

\begin{theorem}[$\epsilon$-$\delta$ and sequential continuity are equivalent]\label{thm:eq}
Let $f : X \to Y$ be a function between metric spaces and let $c \in X$. The
following are equivalent.
\begin{enumerate}
  \item[(A)] $f$ is continuous at $c$ in the sense of Definition~\ref{def:cont}.
  \item[(B)] For every sequence $(x_n) \subseteq X$ with $x_n \to c$, we have $f(x_n) \to f(c)$ in $Y$.
\end{enumerate}
\end{theorem}

\begin{proof}
\noindent\textbf{(A) $\Rightarrow$ (B).}
Assume $f$ is $\epsilon$-$\delta$ continuous at $c$ and let $(x_n)$ be any
sequence with $x_n \to c$. Fix $\epsilon > 0$. By (A), choose $\delta > 0$
such that
\[
  d_X(x, c) < \delta \implies d_Y(f(x), f(c)) < \epsilon. \tag{$\star$}
\]
Since $x_n \to c$, there exists $N \in \mathbb{N}$ with $d_X(x_n, c) < \delta$
for all $n \geq N$. For such $n$, ($\star$) gives
$d_Y(f(x_n), f(c)) < \epsilon$. Hence $f(x_n) \to f(c)$, establishing (B).

\medskip
\noindent\textbf{(B) $\Rightarrow$ (A).}
We argue by contrapositive: suppose $f$ is \emph{not} $\epsilon$-$\delta$
continuous at $c$. Then there exists $\epsilon_0 > 0$ such that for every
$\delta > 0$, some $x \in X$ satisfies $d_X(x, c) < \delta$ yet
$d_Y(f(x), f(c)) \geq \epsilon_0$.

Specialise to $\delta = 1/n$ for each $n \in \mathbb{N}$: pick $x_n \in X$ with
\[
  d_X(x_n, c) < \tfrac{1}{n} \quad \text{and} \quad d_Y(f(x_n), f(c)) \geq \epsilon_0.
\]
The first inequality forces $x_n \to c$. The second forces
$f(x_n) \not\to f(c)$ (otherwise the distances would eventually be smaller
than $\epsilon_0$). This violates (B), establishing the contrapositive and
hence (B) $\Rightarrow$ (A).
\end{proof}

\section{Consequences}

\begin{lemma}[Composition]
If $f : X \to Y$ is continuous at $c$ and $g : Y \to Z$ is continuous at
$f(c)$, then $g \circ f$ is continuous at $c$.
\end{lemma}

\begin{proof}
By Theorem~\ref{thm:eq}, continuity is preserved along sequences: if
$x_n \to c$ then $f(x_n) \to f(c)$, hence $g(f(x_n)) \to g(f(c))$. Apply
Theorem~\ref{thm:eq} in the reverse direction.
\end{proof}

\begin{theorem}[Heine--Cantor, stated]
If $K \subseteq \mathbb{R}^d$ is compact and $f : K \to \mathbb{R}$ is
continuous, then $f$ is uniformly continuous on $K$.
\end{theorem}

The proof (omitted) follows the contrapositive recipe of
Theorem~\ref{thm:eq}: failure of uniformity yields sequences $(x_n), (y_n)$
with $|x_n - y_n| \to 0$ yet $|f(x_n) - f(y_n)| \geq \epsilon_0$; compactness
extracts a convergent subsequence, contradicting pointwise continuity.

\section{Example}

\begin{example}[$x^2$ at $x = 2$]
Fix $\epsilon > 0$. For $|x - 2| < 1$ we have $|x + 2| < 5$, so
\[
  |x^2 - 4| = |x - 2|\,|x + 2| < 5\,|x - 2|.
\]
Choose $\delta = \min(1, \epsilon/5)$. Then
$|x - 2| < \delta \Rightarrow |x^2 - 4| < \epsilon$, verifying
Definition~\ref{def:cont}.
\end{example}

\end{document}
